\section{Data Sources}

This project uses three main kinds of data: tree cadastre records \cite{linzcadastre} \cite{baumkataster2025} Innsbruck \cite{innsbruckcadastre}, \cite{parisarbres} \cite{barcelonaarbrat} \cite{hamburgbaumkataster} , Sentinel-2 satellite imagery \cite{earthenginecodeeditor}, and OpenStreetMap (OSM) urban features \cite{geofabrikeurope}. Together, these sources support the goal of mapping urban tree canopy from freely available data, which were collected, cleaned, and saved in reusable formats across four Jupyter notebooks \cite{scraperbaumkataster} \cite{scrapersentinelosmnotebook} \cite{scrapermoredata}  \cite{scraperdatabase}. The data sources and their cleaning processes are described in the following subchapters.

\subsection{Tree cadastre data}

Tree cadastre datasets from six cities are included: Vienna \cite{baumkataster2025}, Linz \cite{linzcadastre}, Innsbruck \cite{innsbruckcadastre}, Paris \cite{parisarbres}, Barcelona \cite{barcelonaarbrat}, and Hamburg \cite{hamburgbaumkataster}. The datasets were loaded in various formats, including GeoJSON, JSON, CSV, and Excel and provide point-based information about known public trees. They are the main source of tree locations in this project to predict on Sentinel-2 satellite data.\\
\\
Data preparation followed several steps: First, duplicates were checked across all datasets, none were found. Next, columns not needed for modelling were removed, while the following most useful variables were kept: geometry, species name, height, trunk circumference, crown diameter, and a city identifier. Coordinate systems were converted to WGS84 (EPSG:4326) and column names were standardised to make the datasets comparable across cities. Missing values were handled differently depending on the city. Innsbruck, Paris, and Barcelona contain no missing values in the retained columns after cleaning. For Hamburg Port data (hamburg-hpa), crown diameter and circumference contained missing values and rather than dropping these rows, the missing values were imputed using the within-species median, with the overall column median used as a fallback. The same strategy was applied to Linz, where height, crown diameter, and circumference had a small share of missing values, around 4 percent of rows. Missing species names across all datasets were replaced with the placeholder value "unknown".\\
\\
 For Vienna, the source data described missing entries in German as "nicht bekannt", meaning "not known", and "nicht definiert", meaning "not defined". These were also replaced with "unknown". Additionally, young trees recorded with zero height, zero circumference, and zero crown diameter, labelled as ''Jungbaumpflanzung'', were removed, as they represent newly planted trees with no measurable canopy. German species names were stripped from the species column. Vienna reports tree height not as a continuous value but as one of eight categorical height classes. For the exploratory data analysis part, midpoint values were suggested.\\
\\
After cleaning, the cadastre data were saved in GeoPackage and Parquet formats for efficient reuse. Barcelona was saved as Parquet, as it does not have a GeoPandas geometry column.

\subsection{Sentinel-2 satellite imagery}

For Sentinel-2 data collection, first the Copernicus API \cite{copernicus2025} was considered, but then switched to Google Earth Engine \cite{earthenginecodeeditor} to be able to create monthly median composites from Sentinel-2 images filtered by cloud coverage. This makes the process more reproducible and reduces cloud-related problems compared with manual image selection in the previous project. Google Earth Engine also allowed us to restrict the data retrieval to the exact area used in the tree cadastre data. To the six cities mentioned above, five additional cities, Prague, Budapest, Warsaw, Milan, and Athens, were added as new prediction targets for the final database. Their bounding boxes were obtained using the geopy library \cite{geopy}.\\
\\
The project uses seasonal imagery for spring (April), summer (August), and autumn (November), chosen to capture different stages of vegetation development. For Paris and Prague, October was used instead of November, as no suitable November images were available. For Athens, July was used instead of August, as no images were available for that month. A cloud coverage threshold of 20 percent was applied, a stricter 10 percent threshold was tested first but did not produce usable images for several cities and months. From Sentinel-2, four spectral bands are kept: B02 (blue), B03 (green), B04 (red), and B08 (near-infrared), which are the only bands with 10-metre resolution \cite{sentinel2harmonized}. In addition, three vegetation indices are computed: NDVI, EVI, and SAVI \cite{ndvitutorial}. Only 2025 images were used to keep the data consistent with the cadastre sources and the result is a multi-season spectral stack for each city, ready for training and prediction.

\subsection{OpenStreetMap urban features}
Tree cadastre data cover publicly managed trees, so the locations of private trees are unknown. To provide additional spatial context for modelling, urban feature layers were collected from OpenStreetMap (OSM). These data come from Geofabrik downloads \cite{geofabrikeurope} in GeoPackage format. Regional files are available for Austria, Ile-de-France, Catalonia, and Hamburg. From each source, thematic layers were read: roads, buildings, land use, and water.\\
\\
After loading and converting the layers to EPSG:4326 \cite{epsg4326}, they were clipped to the bounding boxes derived from the cleaned cadastre data for each city. This keeps the cadastre data, Sentinel-2 imagery, and OSM layers spatially aligned. At this stage, geometry validity and empty geometries were checked across all layers, no issues were found. No strong filtering of feature classes was applied at this step, in order to preserve flexibility for later modelling. Deeper cleaning is deferred to the exploratory data analysis phase. The cleaned layers were saved as GeoPackage and Parquet files for Vienna, Paris, Hamburg, and Barcelona in one notebook, and for Linz and Innsbruck in a second notebook, using the same Austria GeoPackage file, clipped to the respective city bounds.

\subsection{Additional Viennese Context Data}

Tree cadastre only contains publicly managed trees, while private trees remain unknown and are therefore not available as ground truth. Additionally, OSM data is community-driven and can be spatially incomplete. Since Vienna has particularly good and complete open data provision, three additional official datasets were collected specifically for Vienna to provide richer spatial context for modelling.To help the model learn the distinction between public and potentially private tree locations, the following official datasets were fetched:\\
\\
RNK Land Use Data \cite{landuse} provides official Vienna land-use polygons covering 30 land-use categories grouped into built-up area, green space, and transport. The top-level grouping was used, as it is simpler and more suitable for modelling. German category names were translated to English, and polygon centroids were converted to WGS84.\\
\\
Public Green Space Data \cite{greenspace} contains official polygons of publicly accessible green areas in Vienna. The dataset distinguishes between sub-types such as forest, water, park, playground, dog zone, and sports areas. These sub-type columns were converted to boolean flags, making the dataset more informative for modelling.\\
\\
The Vienna Multi-Purpose Map (Vector Data) \cite{multipurposemap} offers detailed urban surface and infrastructure polygons for Vienna, covering 24 feature types including buildings, roads, green spaces, forests, and water bodies. Duplicate rows, which represented only a small portion of the data, were removed, and rows with missing values were dropped. German feature type names were translated to English.\\
\\
All three datasets were saved as Parquet files, ready for use in exploratory data analysis and modelling.

